\chapter{Camera2D: world coordinates and the Begin/End idea}

\section{The problem with drawing in pixels}

Everything so far has been drawn at screen pixels. That works until the world is
bigger than the window. Then you find yourself writing
\lstinline|DrawRectangle(enemy.x - scrollX, enemy.y - scrollY, ...)| on every
single draw call, and forgetting the subtraction on one of them.

The fix is to stop thinking in pixels. Store everything in \textbf{world}
coordinates --- metres, tiles, whatever your game is made of --- and let a camera
decide which part of the world the window is looking at.

\section{Camera2D}

\begin{lstlisting}
typedef struct Camera2D {
    Vector2 offset;     // where the target lands ON SCREEN (usually the middle)
    Vector2 target;     // the WORLD point the camera is centred on
    float rotation;     // degrees; rotates the whole world
    float zoom;         // 1.0 normal, 2.0 twice as big. Never 0.
} Camera2D;
\end{lstlisting}

The pair that confuses people is \lstinline|target| and \lstinline|offset|. Read
it as a sentence: \emph{``the world point \lstinline|target| is displayed at the
screen position \lstinline|offset|''}. So to centre the camera on the player you
set \lstinline|target| to the player's world position and \lstinline|offset| to
the middle of the window:

\begin{lstlisting}
camera.target = player;                                         // world
camera.offset = (Vector2){ GetScreenWidth()/2.0f, GetScreenHeight()/2.0f };
camera.zoom   = 1.0f;
\end{lstlisting}

\gotcha{\lstinline|zoom| defaults to \lstinline|0| if you write
\lstinline|Camera2D camera = { 0 }| and forget to set it --- and a zoom of zero
scales the entire world down to nothing. Black screen, no error, no clue. If your
Camera2D shows nothing, this is why, every time.}

\section{The Begin/End pair, which is the real lesson}

\begin{lstlisting}
BeginMode2D(camera);
    // WORLD coordinates in here. The camera transforms everything.
    DrawRectangle(player.x, player.y, 30, 30, BLUE);
EndMode2D();

// SCREEN pixels again out here. Nothing below moves with the camera.
DrawText("HEALTH 100", 20, 20, 20, RAYWHITE);
\end{lstlisting}

\keyidea{Inside \lstinline|BeginMode2D|/\lstinline|EndMode2D| you are in world
space and the camera moves things around. Outside, you are in screen pixels and
nothing moves. The HUD goes outside. Always.}

This is worth dwelling on, because in three chapters' time you will meet
\lstinline|BeginMode3D| / \lstinline|EndMode3D| and it is the same idea with a
different camera: the \lstinline|Begin| pushes a transform, the \lstinline|End|
pops it. Your health bar sits outside the pair in both cases, for exactly the
same reason.

You can even nest them in one frame: draw the world in 3D, end the 3D mode, draw
a 2D minimap and the HUD. That is precisely what the mini Doom in chapter 14
does.

\section{Going between the two worlds}

Sooner or later you need to know what the mouse is pointing at in world
coordinates --- to click an enemy, to place a tile:

\begin{lstlisting}
Vector2 worldPoint  = GetScreenToWorld2D(GetMousePosition(), camera);
Vector2 screenPoint = GetWorldToScreen2D(enemy.position, camera);
\end{lstlisting}

The second one is how you put a floating name tag above a character: compute
where the character is on screen, then draw text there \emph{outside} the
Begin/End pair so it keeps a constant size no matter the zoom. The 3D course
uses the 3D version of this same function to label the axes in lesson 1.

\section{Making the camera feel good}

A camera nailed to the player is correct and slightly nauseating. Real games lag
behind a little. One line, using \lstinline|dt| as always:

\begin{lstlisting}
camera.target.x += (player.x - camera.target.x)*5.0f*dt;
camera.target.y += (player.y - camera.target.y)*5.0f*dt;
\end{lstlisting}

Each frame the camera closes a fraction of the gap towards the player. Larger
number, tighter follow. This is called exponential smoothing and it is worth
knowing because it makes almost anything feel better: camera follow, aiming,
health bars catching up to the real value.

\tryit{Lesson 6 gives you a world far bigger than the window, with zoom and
rotation, and a red crosshair drawn \emph{outside} the Begin/End pair so you can
watch it refuse to move with the world.}{./course2d/build.sh 6}
